\chapter{Conclusion \& future work}\label{ch:conclusion}

\section{Conclusion}\label{sec:conclusion}

This thesis asked: \textbf{under which chaos fault classes does pod placement
measurably affect mechanism-level behaviour and user-visible outcomes in a
Kubernetes microservice application, and when do aggregate resilience scores
obscure those effects?} Answered through a pre-registered, Holm-corrected
confirmatory family, the result is a real mechanism whose reach is rigorously
bounded.

The positive thread is the \emph{mechanism}. Placement reproducibly moves a
UDP-conntrack reconvergence signature, and that effect is significant in every
campaign: it is placement-dependent under churn (H2), it interacts with
replication under node-drain (H3's error face, where anti-affine replication
fully rescues the user-route error rate), and it is the one scorecard layer with
near-perfect test-retest reliability (H5, mechanism ICC 0.994 vs a naive
aggregate's 0.066). A NodeLocal DNSCache intervention removes the conntrack drop
under both placements --- the one actionable lever the study confirmed.

The bound is what pre-registration makes precise. No confirmatory hypothesis
clears its registered bar (\S\ref{sec:family-holm}): the dose-response is
detectable but below the 15\,\% SESOI (H1); placement-dependence replicates with
the \emph{opposite} sign, packing concentrating more per-node conntrack churn
than spreading, so the conjunction fails (H2); replication's rescue is clear on
the error face but the trough-depth co-primary could not adjudicate it on this
placement (H3); and the scorecard's availability layer is unreliable under churn
(H5). The latency$\leftrightarrow$availability frontier is likewise
degenerate on the availability face under pod-delete (H4) --- both reflecting
that pod-delete produces no sustained availability outage to discriminate on. Placement moves the
mechanism; at the pre-registered bar, it does not move the user-visible
placement, availability, and dose-response outcomes the hypotheses predicted.

The three contributions of \S\ref{sec:contributions} carry this answer within the
boundary of Chapter~\ref{ch:threats}. The empirical finding --- a real,
DNS-mediated, placement-dependent conntrack mechanism whose user-visible reach is
bounded --- is a statement about a single-replica baseline on one small
virtualized v1.28.6/ipvs cluster (\S\ref{sec:generalizes}), never ``spread is
worse'' or a best strategy. ChaosProbe and the deposit-before-analysis campaign
protocol are the portable parts: any cluster and workload can rerun this design,
and every number traces to an archived, hash-stamped run
(Appendix~\ref{app:provenance}).

The methodological point is the one we would have the field retain. Chaos
evaluation of placement needs \textbf{layered measurement, pre-registration, and
provenance-gated replication, not a single score}: each fault class deposited
its placement signal in a specific layer, a one-number instrument's availability
component was unreliable or undefined exactly where it mattered, and freezing the
hypotheses and effect-size bars before collection is what turned ``we observed
placement effects'' into the more honest ``here is exactly which effects survive
a strict bar, and which do not.''

\section{Future work}\label{sec:future-work}

\subsection*{Second-workload external validity (the highest-value next step)}

The single largest threat (\S\ref{sec:generalizes}) is that all results come from
one workload. The natural next step is to replicate the confirmatory family on a
second microservice application (e.g.\ a Consul/gRPC workload such as
DeathStarBench hotelReservation). A scoping attempt surfaced a concrete blocker
worth recording: that workload's inter-service gRPC clients enter a keepalive
reconnection storm (server-side \texttt{too\_many\_pings} GoAway responses) after
the per-iteration clean-baseline restart, leaving the search path unavailable
for many minutes and false-tainting every iteration's readiness gate, even
though the latency data collected once it recovers is clean. The
restart-for-clean-baseline protocol is thus incompatible with workloads prone to
such storms without one of: a much longer readiness budget, a per-workload skip
of the per-iteration restart (a documented protocol deviation), or tuning the
workload's gRPC keepalive (a change to the system under test). Resolving that
choice is the prerequisite for clean second-workload data.

\subsection*{Decomposing and extending the mechanism}

\begin{itemize}
\item \textbf{Conntrack-drop apportionment.} The drop has both a kernel TCP
  teardown component and the placement-dependent UDP/DNS pool the cache removes
  (\S\ref{sec:h2-mechanism}); a steady-state, multi-iteration protocol-composition
  probe would apportion the two, converting the study's most reproducible signal
  from mechanism-consistent to causally decomposed.
\item \textbf{A rescue-margin the depth face can express.} H3's trough-depth
  co-primary could not adjudicate rescue because the registered absolute one-pod
  margin coincided with the $r{=}1$ dynamic range on this placement; a margin
  defined relative to the realized $r{=}1$ depth (or an integrated outage =
  depth\,$\times$\,duration metric) would let a future pre-registration test the
  depth face on its merits.
\item \textbf{Denser dose levels and other faults.} Finer cross-node-fraction
  steps around the detected (sub-SESOI) trend, and a fault that varies the
  availability face while retaining a latency face, would let the frontier (H4)
  realize the trade-off it is degenerate on here.
\end{itemize}

\subsection*{Environment and scope}

Larger and bare-metal clusters, other CNIs and kube-proxy modes (iptables,
nftables --- the de-scoped direction-transfer comparison), multi-replica
workloads beyond the node-drain regime H3 reaches, production traffic, and
integrating the cross-node fraction into a scheduler as a scoring plugin. The
common thread is the method: each of these runs on the same layered,
pre-registered, provenance-gated protocol this thesis contributes, and each
would sharpen --- not merely extend --- the boundary of where placement matters.
